\documentclass[11pt,a4paper]{article}
\usepackage{curso}
\begin{document}
\portada{Ejercicios --- Lección}{4.4}{Diagnósticos de muestreo}
\section{Ejercicios}
\begin{enumerate}[label=\textbf{E\arabic*.}]
\item (B) Una corrida de $10^5$ barridos tiene $\tau_{\text{int}}=250$ y $\mathrm{Var}(x)=0{,}04$. ESS y barra de error de la media.
\item (B) ¿Por qué Gibbs local se desacelera cerca de $T_c$?
\item (B) Derive la aceptación del intercambio del tempering paralelo a partir del balance detallado en el espacio producto.
\item (B) Escalera geométrica de 8 réplicas de $T=1$ a $T=4$.
\item (I) Para $\rho_k=r^k$, $\tau_{\text{int}}=(1+r)/(1-r)$. ¿Qué $r$ corresponde a $\tau_{\text{int}}=33$?
\item (I) Si un TSU gasta 2 nJ por barrido, ¿cuánta energía cuesta una muestra independiente a $T=3{,}5$ y a $T=2{,}27$ (tabla de los apuntes)?
\item (A) Implemente tempering paralelo para el vidrio $\pm J$ $16\times16$ con 8 réplicas y compare la mejor energía con el recocido de $10^4$ barridos.
\end{enumerate}
\section{Soluciones}
\begin{solucion}
\textbf{E1.} $\mathrm{ESS}=400$; $\sigma_{\bar x}=\sqrt{0{,}04/400}=0{,}01$.

\textbf{E2.} Los dominios abarcan todo el sistema; voltear espines uno a uno los reorganiza con un tiempo que crece como $L^z$.

\textbf{E3.} $\frac{A_{ij}}{A_{ji}}=\frac{e^{-\beta_iE_j-\beta_jE_i}}{e^{-\beta_iE_i-\beta_jE_j}}=e^{(\beta_i-\beta_j)(E_i-E_j)}$.

\textbf{E4.} $T_k=4^{k/7}$: 1; 1,22; 1,49; 1,81; 2,21; 2,69; 3,28; 4.

\textbf{E5.} $r=(\tau-1)/(\tau+1)=32/34=0{,}941$.

\textbf{E6.} $2\times8{,}5=17$ nJ; $2\times2200=4{,}4\ \mu$J.

\textbf{E7.} El tempering suele alcanzar energías $\le-1{,}37$ con un costo comparable; reportar la mejor energía y la fracción de réplicas que la encuentran.
\end{solucion}
\end{document}
